% =====================================================================
\chapter{RMSProp}
\label{ch:rmsprop}
% =====================================================================

\tikzset{bwcontours/.pic={
  \foreach \r in {0.5,1.05,1.65} \draw[very thick] (0.4,0) ellipse ({2.3*\r} and {\r});
  \fill[inkred] (0.4,0) circle (2.5pt);
  \draw[inkgreen,thick,-Stealth] (-3.2,-1.45) -- (-3.2,1.6) node[above,font=\small] {$b$};
  \draw[inkgreen,thick,-Stealth] (-3.2,-1.45) -- (1.5,-1.45) node[right,font=\small] {$w$};}}

\lettrine[lines=2,findent=2pt,nindent=0pt]{\color{inkred}A}{s} we have seen before with momentum, with the EMA we turned a noisy pattern into a smooth one. RMSProp relies on the same idea, but with only one important difference.\ednote{EMA (exponential moving average) and EWA (exponentially weighted average) are two names for the same thing; the notes use both.}

The idea of RMSProp comes from this figure:
\begin{center}
\begin{tikzpicture}
  \pic at (0,0) {bwcontours};
  \draw[inkblue,thick,-Stealth] (-2.5,-0.2) -- (-2.2,1.0) -- (-1.8,-0.9) -- (-1.4,1.0) -- (-1.0,-0.85) -- (-0.6,0.9) -- (-0.25,-0.6) -- (0.05,0.55) -- (0.35,0.05);
\end{tikzpicture}
\end{center}

So the remark here is that, at each step, the gradient is bigger towards $b$, and less important towards $w$. We actually want it to go in the direction of $w$! So to solve this, we introduce the RMSProp algorithm:

\Def[Normalise each direction by its historical volatility.]{Def: RMSProp}

RMSProp stands for \textbf{R}oot \textbf{M}ean \textbf{S}quare \textbf{Prop}agation, and this is the idea:

\begin{center}
\begin{tikzpicture}
  \node (m) at (0,0) {\Large $v_t=\beta v_{t-1}+(1-\beta)\,\nabla f_t^{\,2}(x_t)$};
  \node[inkgold,font=\small] (lm) at (-4.6,0) {Mean $\rightarrow$};
  \node[inkgold,font=\small] (sq) at (2.3,1.0) {Square};
  \draw[-Stealth,inkgold] (sq) -- (1.85,0.35);
  \node[inkred,font=\small,align=left,anchor=west] at (3.4,0.2) {\warn\ Pointwise squaring!\\ \phantom{\warn\ }$\nabla f_t(x_t)\odot\nabla f_t(x_t)$};
  \node (u) at (0,-1.8) {\Large $x_{t+1}=x_t-\gamma\,\dfrac{\nabla f_t(x_t)}{\sqrt{v_t}+\varepsilon}$};
  \node[inkred,font=\small,align=center] (pr) at (-3.6,-3.2) {\warn\ ``Pointwise\\ root''};
  \draw[-Stealth,inkred] (pr.east) -- (0.65,-2.35);
  \node[inkgold,font=\small] (rt) at (2.9,-3.1) {Root};
  \draw[-Stealth,inkgold] (rt.west) -- (0.9,-2.45);
\end{tikzpicture}
\end{center}

\Q{Well, what does this mean, and how would that help us reduce the effect of large gradients towards $b$?}

Let's see the details:
\[
v_t=\beta v_{t-1}+(1-\beta)\,\nabla f_t^{\,2}(x_t)
\]
\[
\swarrow\hspace{4cm}\searrow
\]
\[
v_{t,b}=\beta v_{t-1,b}+(1-\beta)\left(\frac{\partial f}{\partial b}\right)^2
\qquad\Big|\qquad
v_{t,w}=\beta v_{t-1,w}+(1-\beta)\left(\frac{\partial f}{\partial w}\right)^2
\]

\hl{If the step is bigger towards $b$:}
\[
\left(\frac{\partial f}{\partial b}\right)^2>\left(\frac{\partial f}{\partial w}\right)^2
\implies v_{t,b}>v_{t,w}\implies\sqrt{v_{t,b}}>\sqrt{v_{t,w}}
\]

And to compute the step:
\begin{center}
\begin{tikzpicture}
  \node at (0,0) {$\left\{\begin{aligned}
    x_{t+1,b}&=x_{t,b}-\gamma\,\dfrac{\partial f/\partial b}{\sqrt{v_{t,b}}+\varepsilon}\\[4pt]
    x_{t+1,w}&=x_{t,w}-\gamma\,\dfrac{\partial f/\partial w}{\sqrt{v_{t,w}}+\varepsilon}
  \end{aligned}\right.$};
  \node[inkgreen,font=\small\itshape,align=left,anchor=west] (a) at (3.2,0) {We divided by something\\ bigger than this one.};
  \draw[-Stealth,inkgreen] (a.west) to[out=160,in=0] (2.3,0.55);
  \draw[-Stealth,inkgreen] (a.west) to[out=200,in=0] (2.3,-0.6);
\end{tikzpicture}
\end{center}

So if the gradient is big in a direction, we divide by a big quantity to reduce it, and that leads to this:

\begin{center}
\begin{tikzpicture}
  \pic at (0,0) {bwcontours};
  \draw[inkblue,thick] (-2.5,-0.2) -- (-2.2,1.0) -- (-1.8,-0.9) -- (-1.4,1.0) -- (-1.0,-0.85) -- (-0.6,0.9) -- (-0.25,-0.6) -- (0.05,0.55) -- (0.35,0.05);
  \draw[inkpurple,very thick,-Stealth] (-2.5,-0.2) -- (-1.9,0.15) -- (-1.3,-0.1) -- (-0.7,0.1) -- (-0.1,-0.05) -- (0.3,0.0);
  \node[inkgold,font=\small\itshape,anchor=west] at (4.6,0.4) {$\leftarrow$ Propagation};
  \node[font=\small] at (0,-2.1) {{\color{inkblue}$\bm{\wr\wr}$} : SGD \qquad {\color{inkpurple}$\bm{\wr\wr}$} : RMSProp};
\end{tikzpicture}
\end{center}

\Q{Another thing: what's ``$\varepsilon$''?}

Epsilon plays a crucial role. It has two big roles:
\begin{enumerate}[label=\arabic*/]
  \item \textbf{Avoid dividing by zero!} (In practice we take $\varepsilon=10^{-8}$; PyTorch: $\varepsilon=10^{-8}$, TensorFlow: $\varepsilon=10^{-7}$.)
  \item \textbf{Defining a minimum of division.} Imagine $\sqrt{v_{t}}\simeq10^{-20}$.
  \begin{itemize}
    \item Without $\varepsilon$: $\dfrac{\nabla f_t(x_t)}{10^{-20}}=\nabla f_t(x_t)\times10^{20}$ \quad ``\hlr{Explosion}''.
    \item With $\varepsilon$: $\dfrac{\nabla f_t(x_t)}{\varepsilon+10^{-20}}\simeq\dfrac{\nabla f_t(x_t)}{\varepsilon}$ \quad (if $\varepsilon=10^{-8}$: $10^{-8}+10^{-20}\simeq10^{-8}$).
  \end{itemize}
\end{enumerate}
Those are the main roles of $\varepsilon$.

\Q{One last question: why do we divide by $\sqrt{v_t}$ and not $v_t$?}

And the answer comes from maths. With $v_0=0$:\ednote{The notes write the sum up to $+\infty$ and then use the finite sum $\frac{1-\beta^t}{1-\beta}$. Starting from $v_0=0$, $v_t$ has exactly $t$ terms, from $k=0$ to $t-1$.}
\begin{align*}
v_t&=\beta v_{t-1}+(1-\beta)\,\nabla f_t^{\,2}(x_t)\\
\implies v_t&=(1-\beta)\sum_{k=0}^{t-1}\beta^k\,\nabla f_{t-k}^{\,2}(x_{t-k})\\
E(v_t)&=(1-\beta)\sum_{k=0}^{t-1}\beta^k\,E\bigl(\nabla f_{t-k}^{\,2}(x_{t-k})\bigr)
\end{align*}
Here comes the assumption of optimization in deep learning; we assume that:
\[
E\bigl(\nabla f_t^{\,2}(x_t)\bigr)\simeq E\bigl(\nabla f_{t-1}^{\,2}(x_{t-1})\bigr)\simeq E\bigl(\nabla f^{\,2}(x)\bigr)
\]
That leads to:
\begin{align*}
E(v_t)&\simeq(1-\beta)\times\frac{1-\beta^t}{1-\beta}\,E\bigl(\nabla f^{\,2}(x)\bigr)\\
\beta<1&\implies\lim_{t\to+\infty}\beta^t=0\\
&\boxed{E(v_t)\simeq E\bigl(\nabla f^{\,2}(x)\bigr)}
\end{align*}

And for the $w$ direction:
\[
E(v_{t,w})\simeq E\left(\left(\frac{\partial f}{\partial w}\right)^2\right)
\qquad\implies\qquad
[v_{t,w}]\simeq\frac{[f]^2}{[w]^2}
\]
(where $[\,\cdot\,]$ is the unit, the ``dimension'').

So, let's say we don't compute the root:
\begin{center}
\begin{tikzpicture}
  \node (e) at (0,0) {$[x_{t+1,w}]=[x_{t,w}]-\gamma\,\dfrac{\;\dfrac{[f]}{[w]}\;}{\dfrac{[f]^2}{[w]^2}+\varepsilon}$};
  \node[inkpurple,font=\small\itshape,anchor=west] (d) at (3.4,0) {Dimension problem};
  \draw[-Stealth,inkpurple] (d.west) to[out=160,in=0] (2.0,0.45);
  \draw[-Stealth,inkpurple] (d.west) to[out=200,in=0] (2.3,-0.4);
\end{tikzpicture}
\end{center}
But with the root, $[\sqrt{v_{t,w}}]=\dfrac{[f]}{[w]}$:
\begin{center}
\begin{tikzpicture}
  \node (e) at (0,0) {$[x_{t+1,w}]=[x_{t,w}]-\gamma\,\dfrac{\;\dfrac{[f]}{[w]}\;}{\dfrac{[f]}{[w]}+\varepsilon}$};
  \node[inkpurple,font=\small\itshape,anchor=west] (d) at (3.3,0) {Fixed!};
  \draw[-Stealth,inkpurple] (d.west) to[out=160,in=0] (2.0,0.45);
  \draw[-Stealth,inkpurple] (d.west) to[out=200,in=0] (2.2,-0.4);
\end{tikzpicture}
\end{center}
The ratio has no unit anymore: the step only depends on $\gamma$, whatever the scale of $f$ or of $w$.\ednote{Sentence added to make the conclusion of the dimensional analysis explicit.}
